\chapter{2025:Susumu Kitagawa and Richard Robson and Omar M. Yaghi}

\label{chap:chemistry2025}

The Nobel Prize in Chemistry for the year 2025 was awarded jointly to Susumu Kitagawa and Richard Robson and Omar M. Yaghi.

\textbf{Motivation:}

for the development of metal–organic frameworks

\section{Susumu Kitagawa}

\subsection*{Early Life and Education}

Susumu Kitagawa was born in 1951 in Kyoto, Japan.

\subsection*{Career and Major Research}

Kitagawa studied chemistry at Kyoto University and became a professor there. He pioneered porous coordination polymers, the family of materials now known as metal-organic frameworks, and showed that their cavities can reversibly take up and release gases and other small molecules.

\subsection*{Nobel Prize-winning Work}

The work for which Susumu Kitagawa was awarded the Nobel Prize in Chemistry in 2025 was described by the Nobel Committee as: "for the development of metal–organic frameworks".

Kitagawa demonstrated that coordination polymers can form crystalline frameworks with permanent porosity, and his studies established how their flexible structures respond to guests such as gas molecules.

\subsection*{Significance and Impact}

His work on porous coordination polymers revealed that metal-organic frameworks could be tuned for applications such as gas storage and separation.

\subsection*{Later Career and Honors}

Kitagawa is professor emeritus of Kyoto University, where he continues research on porous materials.

\subsection*{Legacy}

Kitagawa's concept of porous coordination polymers is a foundation of metal-organic framework chemistry and of modern porous materials science.

\subsection*{Selected Publications}

"Functional Porous Coordination Polymers" (Angewandte Chemie International Edition, 2004)

\clearpage

\section{Richard Robson}

\subsection*{Early Life and Education}

Richard Robson was born in 1937 in Glusburn, England.

\subsection*{Career and Major Research}

Robson earned bachelor's and doctoral degrees from the University of Oxford, receiving his doctorate in 1962. He became a professor of inorganic chemistry at the University of Melbourne, where in the late 1980s he designed coordination polymer networks with predictable topologies by linking metal centers with rigid organic units.

\subsection*{Nobel Prize-winning Work}

The work for which Richard Robson was awarded the Nobel Prize in Chemistry in 2025 was described by the Nobel Committee as: "for the development of metal–organic frameworks".

Robson showed how to assemble metal ions and organic linkers into crystalline networks with large channels, establishing the design principles for what became metal-organic frameworks.

\subsection*{Significance and Impact}

His networks demonstrated that the architecture of a porous solid could be planned in advance from its building blocks, a central idea of the field.

\subsection*{Later Career and Honors}

Robson was professor of inorganic chemistry at the University of Melbourne until his retirement in 2004.

\subsection*{Legacy}

Robson's design rules for infinite framework structures are a cornerstone of the chemistry of coordination polymers and metal-organic frameworks.

\subsection*{Selected Publications}

"Assembly of Porphyrin Building Blocks into Network Structures with Large Channels" (Nature, 1994)

\clearpage

\section{Omar M. Yaghi}

\subsection*{Early Life and Education}

Omar M. Yaghi was born in 1965 in Amman, Jordan.

\subsection*{Career and Major Research}

Yaghi received his doctorate from the University of Illinois at Urbana-Champaign in 1990. He held professorships at Arizona State University, the University of Michigan, and the University of California, Los Angeles, before becoming a professor at the University of California, Berkeley. He synthesized the highly porous metal-organic framework MOF-5 in 1999 and founded the field of reticular chemistry.

\subsection*{Nobel Prize-winning Work}

The work for which Omar M. Yaghi was awarded the Nobel Prize in Chemistry in 2025 was described by the Nobel Committee as: "for the development of metal–organic frameworks".

Yaghi's metal-organic frameworks combine metal clusters and organic linkers into extended porous crystals with record surface areas, and his reticular chemistry approach makes their structures predictable and tunable.

\subsection*{Significance and Impact}

His robust, high-surface-area frameworks are used for storing gases such as hydrogen and methane, capturing carbon dioxide, and harvesting water from desert air.

\subsection*{Later Career and Honors}

Yaghi continues his research on reticular chemistry at the University of California, Berkeley.

\subsection*{Legacy}

The metal-organic frameworks Yaghi developed are among the most porous materials known and are being commercialized for energy and environmental applications.

\subsection*{Selected Publications}

"Design and Synthesis of an Exceptionally Stable and Highly Porous Metal-Organic Framework" (Nature, 1999)

\clearpage

\section*{Chapter Summary}

The 2025 prize recognized the development of metal-organic frameworks, porous crystals built from metal ions and organic linkers. Kitagawa pioneered porous coordination polymers, Robson established the design of network structures, and Yaghi synthesized high-porosity frameworks with predictable architectures, materials now used for gas storage, carbon capture, and catalysis.
